%% 
%% Dissertation
%% 


\documentclass[a4paper,12pt,twoside,openright]{report}


%%
%% Definitions of names and suchlike
%%

\def\authorname{David J.\ Barker\xspace}
\def\authorcollege{Jesus College\xspace}
\def\authoremail{David.Barker@cl.cam.ac.uk}
\def\dissertationtitle{Developing a formally verified algorithm for static register allocation}
\def\wordcount{TBC}


\usepackage{epsfig,graphicx,parskip,setspace,tabularx,xspace,holtexbasic,alltt}

%% START OF DOCUMENT
\begin{document}


%% FRONTMATTER (TITLE PAGE, DECLARATION, ABSTRACT, ETC)
\pagestyle{empty}
\singlespacing
% title page information
\begin{titlepage} 

\begin{center}
\noindent
\huge
\dissertationtitle \\
\vspace*{\stretch{1}}
\end{center}

\begin{center}
\noindent
\huge
\authorname \\
\Large
\authorcollege      \\[24pt]
%\includegraphics{CUni3.eps}
\end{center}

\vspace{24pt}

\begin{center}
\noindent
\large
{\it A dissertation submitted to the University of \\ 
Cambridge in partial fulfilment of the requirements \\
for Part III of the Computer Science Tripos} 
\vspace*{\stretch{1}}
\end{center}

\begin{center}
\noindent
University of Cambridge \\
Computer Laboratory     \\
William Gates Building  \\
15 JJ Thomson Avenue    \\
Cambridge CB3 0FD       \\
{\sc United Kingdom}    \\
\end{center}

\begin{center}
\noindent
Email: \authoremail \\
\end{center}

\begin{center}
\noindent
\today
\end{center}

\end{titlepage} 

%\newpage
%\vspace*{\fill}

\onehalfspacing
\newpage
{\Huge \bf Declaration}

\vspace{24pt} 

I \authorname of \authorcollege , being a candidate for the Part III in Computer
Science, hereby declare that this report and the work described in it
are my own work, unaided except as may be specified below, and that
the report does not contain material that has already been used to
any substantial extent for a comparable purpose.

\vspace{24pt}
Total word count: \wordcount

\vspace{60pt}
\textbf{Signed}: 

\vspace{12pt}
\textbf{Date}:


\vfill

This dissertation is copyright \copyright 2014 \authorname. 
\\
All trademarks used in this dissertation are hereby acknowledged.



%\newpage
%\vspace*{\fill}

\singlespacing
\newpage
{\Huge \bf Abstract}
\vspace{24pt} 


This dissertation describes the implementation and verification of a full register allocation algorithm, designed to be representative of the sort of allocator that might be used in a practical optimising compiler. A model of three-address code (and its execution behaviour) was developed and a verified live variable analysis function and clash graph generator were constructed. A number of different graph colouring algorithms were then implemented and proved correct, as well as various heuristics. The work was then tied together by an overarching proof of correctness which demonstrates that the algorithm does not change code behaviour. The algorithm was also extended to construct preference graphs based on `move' instructions and use them in colouring, and to handle spilling of registers where not enough physical registers are available, and these additions were also verified in depth.


\newpage
\vspace*{\fill}


\pagenumbering{roman}
\setcounter{page}{0}
\pagestyle{plain}
\tableofcontents
\listoffigures
\listoftables

\onehalfspacing

%% START OF MAIN TEXT 

\chapter{Introduction}
\pagenumbering{arabic}
\setcounter{page}{1}

This is the introduction where you should introduce your work.  In
general the thing to aim for here is to describe a little bit of the
context for your work --- why did you do it (motivation), what was the
hoped-for outcome (aims) --- as well as trying to give a brief
overview of what you actually did.

It's often useful to bring forward some ``highlights'' into 
this chapter (e.g.\ some particularly compelling results, or 
a particularly interesting finding). 

It's also traditional to give an outline of the rest of the
document, although without care this can appear formulaic 
and tedious. Your call. 


\chapter{Background} 

\section {The problem of register allocation}

When compiling code, we generally assume the existence of an infinite number of
locations in which to store values. Each variable used by the program is treated
as its own location, and temporaries and results are freely allocated to conceptual
registers. However, in reality all processors have only finite numbers of physical
registers, some of which are reserved for specific purposes, and so the compiler
ultimately has to map variables in the code to registers in such a way that ensures
the code's behaviour is the same.

There are many issues involved when doing this. The main problem is that code will
often use far more `virtual' registers than actually exist on the target processor.
The simplest solution to this is to store any values which won't fit in memory, but
this considerably worsens performance. The solution is to find variables which are
never in use at the same time and place them in the same register, allowing registers
to be re-used when the values originally assigned to them are not in use.

This optimisation is enabled by liveness analysis. Essentially, a variable is live
where changing its value will affect the program's future behaviour [ref?]. Because
this is difficult to decide with absolute precision, we generally compute a safe
over-approximation of liveness: a variable is live if an instruction reachable from
here reads its value, and its value is not written to in a preceding instruction.
Liveness is computed by repeatedly iterating over instructions in the code from the
end backwards, updating each instruction's set of live variables according the
following liveness equation:

\begin{center}
$live(n) = \left(live(n+1) \setminus write(n)\right) \cup read(n)$
\end{center}

(Where live(n) is the set of variables live at instruction n, live(n+1) is those
live at the next instruction and read(n) and write(n) are the variables read and written
by the instruction.)

This is repeated until the live variable sets stop changing. With this information,
we can deduce which variables are live at the same time, and thus which variables
cannot occupy the same register because doing so would result in one overwriting the
other whilst the other was about to be used.

When no allocation can satisfy these clashes between variables, it becomes necessary
to spill variables to memory. This means reading or writing said variables takes
considerably longer, and so an effective register allocation algorithm should minimise
spills wherever possible.

%TODO: [Find a reference or two for this next bit?]
Another thing to consider is that some target architectures have constraints on what
registers can be used for. For instance, a particular operator may require its operands
to be stored in particular registers, or may place the output in a designated `result'
register. This means move instructions will be required to satisfy these requirements,
slowing the overall program down. Good register allocators allow for this by biasing
allocation so that these move instructions can be eliminated wherever possible.

\section{Isomorphism to the graph-colouring problem}

The most common approach to register allocation is to reduce the problem to one of graph
colouring, where we have a graph which we wish to colour with a maximum of K colours such
that no two adjacent vertices share the same colour (TODO ref). Given the knowledge of which
variables are live at the same time, we construct a clash graph in which the vertices
represent registers and an edge between two vertices exists where the corresponding
registers are simultaneously live. Treating a colour as representing a physical register
on the target machine, we then attempt to `colour' the graph with target registers such
that vertices which are linked by a clash edge are always given different colours.

A successful colouring will naturally give a mapping from virtual to physical registers
enforcing that no registers which are simultaneously live are placed in the same physical
register. Formally verifying that the colouring resulting from this process will never
affect the behaviour of the code being compiled was the core objective of this project.

The K-colourability problem is known to be NP-complete wherever K is greater than two,
and so finding a truly optimal solution is rarely feasible in practice. Instead,
register allocation algorithms are designed primarily to find a solution which is
correct, and from there use heuristics to find a solution which is as optimal as
can be achieved in reasonable time. There are several approaches to this, and many
of them were verified in detail in this project. As discussed later, the definitions
and proofs in the project are designed to be largely independent of the algorithm
and so extending it to verify some of the algorithms not covered would be relatively
simple.

\section {Basic approaches to colouring}

The simplest way of colouring a graph is to take each vertex in turn and assign
it a new colour until all vertices have been coloured. This approach is guaranteed
to generate a valid colouring (as verified in Section (TODO section)), but is also very
inefficient as registers are never re-used for variables which could reside in the
same register. In a system with few physical registers this will result in registers
being spilled to memory, and as the order colours are assigned is arbitrary this could
result in frequently-used registers being stored in memory causing constant loads and
stores and poor performance.

An improvement is to take each vertex in turn and assign it the lowest register which
isn't already in use by a vertex with which it clashes. This is fairly simple to
implement, and significantly improves register utilisation as it re-uses old registers
whenever the opportunity presents itself (TODO ref?).

However, lowest-first allocation doesn't always find opportunities to re-use old
registers. A further refinement is to check whether we can avoid using a new register
by swapping colours of already-allocated vertices. This often improves register
utilisation further, as sub-optimal decisions made early on in the process can be
changed when colouring new vertices later.

The algorithm resulting from applying all these refinements will work fairly well,
but in practice the quality of the colouring is often also heavily dependent on the
order in which we colour vertices. Decisions made at the beginning naturally constrain
decisions made later, even if we allow the algorithm to swap existing colours. We
therefore need to use a heuristic ordering step at the beginning to improve the order
in which vertices are presented to the main colouring algorithm. Some such heuristics
are described in the next section.

\section{Heuristics}
Heuristics are used in graph-colouring algorithms to select a vertex ordering which
gives a more optimal solution. There are several heuristics which are used in practice,
some of which were formally verified during the course of this project.

One of the simplest heuristics is to order the vertices by their degree. By presenting
the vertices to the colouring algorithm in descending order of degree, vertices which
are most likely to clash with others are coloured first so that those left over towards
the end will be easier to colour without assigning new colours. There are also a few
variations on this: saturation degree ordering, for instance, picks the next vertex to
colour based on which vertex is currently adjacent to the most different colours.
Incidence degree colouring is another variant which sorts based on the number of adjacent
vertices which have already been coloured (TODO ref).

A more complex heuristic, which is useful when attempting to colour with at most K colours,
is to pick at each stage the vertex of highest degree less than K (this constraint on its
degree ensures that it can definitely be coloured if only K colours are available). This
vertex is then removed from the graph and placed on top of a vertex stack. The process
repeats, each time removing a vertex which is guaranteed to be colourable, until no
vertices of degree less than K remain. Those which have not yet been removed are spilled
to memory, and the rest of the vertices are popped off the stack one by one and passed
to the colouring algorithm. The heuristic essentially priorities registers with few
conflicts, spilling those with many conflicts to memory, and colours those of lowest
degree first.

This is not an exhaustive list of the heuristics used in practice, but these two types
are very common and so verification focussed primarily on these. As mentioned, most of
the verification conditions were designed to be as general as possible such that other
more exotic heuristics could also be verified without much additional effort.



\chapter{Related Work} 

TODO


\chapter{Live variable analysis and clash graph generation}

\section*{Introduction}

In order to achieve a complete end-to-end verification, it was decided that the verified
allocator should take in some representation of three-address code (typical of this stage
in the compilation process), and output three-address code with registers allocated. The
overall proof would therefore show that the behaviour of the resulting code was exactly
the same as that of the source in all situations. For this reason, a simplified
representation of three-address code and an evaluation function were developed and an
algorithm to generate clash graphs was verified. This chapter explores the representation
used, the way `colouring' registers was represented and the development and verification
of the clash graph generator, along with the way a `correct' register allocation was defined.
The proofs showing that these definitions are correct and linking them to the project�s overall
goal will follow in Chapter (TODO: bringing everything together).

\section{Three-address code representation}

Here we explore the simplified representation of three-address code, how evaluation was
defined and how colourings were modelled and applied.

\subsection{Registers and instructions}

In real systems, register naming conventions tend to vary between processors, and there
are often multiple groups of registers available. For the purposes of verification, this
was simplified by using natural numbers to refer to registers, as in practice these numbers
could trivially be mapped to real register names. Furthermore, there will generally be a
great many different operations available on a given machine, but as the focus is on register
allocation the only thing that matters here is the registers read and those assigned.
Compilers commonly use three-address code as an intermediate representation, so a given
instruction essentially consists of two source registers and one destination register. This
led to the very simplified instruction representation given below:

\begin{alltt}\small
	\HOLTyOp{inst} = \HOLConst{Inst} \HOLKeyword{of} \HOLTyOp{num} \HOLTokenImp{} \HOLTyOp{num} \HOLTokenImp{} \HOLTyOp{num}
\end{alltt}

For the main verification task, this was entirely satisfactory as it captures all the
information needed. Later work extended this representation slightly: for instance, the
addition of preferences meant it was necessary to include move instructions so that these
could be used to compute preference graphs. Extensions to the three-address representation
are discussed in Chapter (TODO: extensions chapter).

Code blocks were then represented as simple lists of these instructions, to simplify
verification. The extension to more complex code featuring basic blocks is not too great
a leap, and is discussed in Section (TODO: extensions-complex code). The main difference
in using this representation is that live variable analysis is somewhat simpler.

\subsection{Code evaluation}

As mentioned, it was not necessary for verification to have a wide range of instructions.
The register allocation can be proven correct simply by showing that evaluating each
instruction with an arbitrary binary operator and storing the result in the destination
register has the same effect as it did before allocation. The precise meaning of `has
the same effect' is explained in Section (TODO bringing together - proving stuff about code
evaluation), where it is proven to hold for the register allocator built during this project.
The evaluation function used is given below:

\begin{alltt}\small
	\HOLConst{eval} \HOLFreeVar{f} \HOLFreeVar{s} [] = \HOLFreeVar{s}
\HOLConst{eval} \HOLFreeVar{f} \HOLFreeVar{s} (\HOLConst{Inst} \HOLFreeVar{w} \HOLFreeVar{r\sb{\mathrm{1}}} \HOLFreeVar{r\sb{\mathrm{2}}}::\HOLFreeVar{code}) =
\HOLConst{eval} \HOLFreeVar{f} ((\HOLFreeVar{w} =+ \HOLFreeVar{f} (\HOLFreeVar{s} \HOLFreeVar{r\sb{\mathrm{1}}}) (\HOLFreeVar{s} \HOLFreeVar{r\sb{\mathrm{2}}})) \HOLFreeVar{s}) \HOLFreeVar{code}
\end{alltt}

Here `f' is an arbitrary binary operator and `s' is the store, a total function representing
the values of all registers. Each instruction is evaluated by looking up the two source
registers in the store and applying the binary operation to them, then updating the store
function to map the destination register to the resulting value. The base case, where we
have reached the end of the code block, simply returns the store as it is.


\chapter{Colouring algorithms}

TODO


\chapter{Bringing everything together}

TODO


\chapter{Extension features}

TODO


\chapter{Summary and conclusions}

TODO




\appendix
\singlespacing

\bibliographystyle{unsrt} 
%\bibliography{dissertation} 

\end{document}
